\chapter{超局所台と超局所化}
\newcommand{\mucat}[2][]{\mathsf{D}^{\mathrm{b}}_{#1}(#2)}

\section*{要約}
\section{圏\(\Dompb(X;\varOmega)\)}
\(V\)を\(\cotb[X]\)の部分集合とする．
\(\Dompb(X)\)の充満部分圏\(\mucat[V]{X}\)を次で定める．
\begin{equation}
    \Ob(\mucat[V]{X})\coloneqq
    \left\{F\in\Ob(\Dompb(X));\muS(F)\subset V\right\}
\end{equation}
\(\mucat[V]{X}\)は，
\(\mres[{[1]}]{\mucat[V]{X}}\)をシフト関手，
\[
    \left\{F\to G\to H\to F[1]\mathrel{;}\text{d.t. in \(\Dompb(X)\)}, F,G,H\in\mucat[V]{X}\right\}
\]
を完全三角の族として，三角圏になる．
\section{余接束内の法錐}
\section{順像}
\section{超局所化}
\section{包合性と伝播}
\section{包合的多様体の近傍における層}
\section{超局所化と逆像}